\section{Schemat przeprowadzonych eksperymentów}

\subsection{Podział danych}

W eksperymentach wykorzystano zbiór Wine, który został podzielony
stratyfikowanie na trzy części:
\begin{itemize}
    \item zbiór treningowy,
    \item zbiór walidacyjny,
    \item zbiór testowy.
\end{itemize}

\begin{table}[H]
    \centering
    \begin{tabular}{lc}
\toprule
Podzbiór & Udział w danych \\
\midrule
Treningowy & 60\% \\
Walidacyjny & 20\% \\
Testowy & 20\% \\
\bottomrule
\end{tabular}

    \caption{Schemat podziału danych wykorzystywany w etapie trzecim}
    \label{tab:data_split_summary}
\end{table}

\subsection{Schemat porównania konfiguracji}

Dla każdej analizowanej konfiguracji wykonywany jest osobny proces uczenia.
W trakcie treningu zapisywane są:
\begin{itemize}
    \item wartości funkcji kosztu dla kolejnych epok,
    \item metryki jakości na zbiorze treningowym,
    \item metryki jakości na zbiorze walidacyjnym,
    \item metryki jakości na zbiorze testowym.
\end{itemize}

\subsection{Metryki oceny}

Do oceny jakości modeli wykorzystano następujące metryki:
\begin{itemize}
    \item wartość funkcji kosztu,
    \item accuracy,
    \item sensitivity liczone osobno dla każdej klasy,
    \item specificity liczone osobno dla każdej klasy,
    \item wartości uśrednione \textit{macro sensitivity} i
    \textit{macro specificity}.
\end{itemize}

\subsection{Artefakty zapisywane w trakcie eksperymentów}

Dla każdej konfiguracji zapisywane są między innymi:
\begin{itemize}
    \item przebieg funkcji kosztu w kolejnych epokach,
    \item pliki CSV z historią treningu,
    \item zbiorcze zestawienia wyników końcowych,
    \item zestawienia średnich i odchyleń standardowych pomiędzy uruchomieniami.
\end{itemize}
